\documentclass[phy1200]{handout}

\title{Lorentz transformations and velocity addition}

\begin{document}

\maketitle

One popular way to visualize the reference frames talked about in special relativity is using spacetime diagrams, also known as Minkowski diagrams.\footnote{Minkowski was one of Einstein's professors in undergraduate, and an early proponent of Einsteins theories.} Usually we only consider one spatial dimension + time in such diagrams as 3D drawing is hard. Position $x$ is on the x-axis, and time $t$ is on the y-axis. Additionally, for these diagrams, we usually assume $c=1$.\footnote{Sometimes instead of a time axis some people draw a $ct$ axis, which has the same effect.} On schoology, you will find a desmos graph where you can play around with Minkowski diagrams

\section{Reading a Minkowski diagram}
Points on the diagram are \textbf{events}, happening at a specific point in space and specific time.

As time proceeds, an object's position traces out a \textbf{world line} (basically a position-time graph turned on its side). Make a world line in the Desmos tool in the apporpriate folder. Consider what the following world lines would look like:
\begin{enumerate}
    \item A stationary object at $x=0$.
    \item A stationary object at some other location.
    \item An object moving with some velocity $v<c$.
    \item A beam of light traveling at $v=c$.\mgpar{World lines moving at the speed of light are called light cones.}
\end{enumerate}
Given that things can't travel faster than the speed of light, what world lines would be impossible?

Another set of lines on the diagram are simultaneity lines. Where are all the events that take place at a given time (start with $t=0$ and move up to general $t$).

\section{Lorentz Transformations}
The Lorentz transformations encode the length contraction and time dilation we discussed in class, but it applies it to all possible ticks of the moving clock (the $t'$ axis) and all possible measuring rods passing the origin (the $x'$ axis).

Lets denote by $(x,t)$ a location in the stationary frame (based on stationary measuring rods and stationary clock ticks). We will denote by $(x',t')$  a location in the moving frame (based on measuring rods and clocks moving at speed $v$ relative to the stationary frame). These coordinate frames are related by the Lorentz transformation:
\begin{align*}
    x'&=\gamma (x-vt)\\
t'&=\gamma \left(ct- \frac{vx}{c^2}\right)
\end{align*}
What are the inverse Lorentz transformations (solve this system of equations for $x,t$)? Why does this make sense in terms of the principle of relativity?

We could make a new grid of all the ticks of our moving clock ($t'=0,1,2,\dots$) and how far points are according to the moving measuring rod ($x'=0,1,2,\dots$). These will no longer be straight up and down. 

To explore the stationary frame representation of the moving grid, we would want formulas for $x(x',t)$ so that we can put in the measuring rod values $x'$ and still know what it looks like on our graph. Similarly, $t(t',x)$. Put these on lines 4-5.

Turn off the rest frame grid and turn on the the moving grid. Explain what the $x'=const.$ lines look like in the stationary frame (they are world lines corresponding to). Why does this make sense?

Notice that the set of all events happening at the same stationary time $t=0$ and the set of events happening at the same moving time are not the same. Explain the term `relativity of simultaneity'.

\section{Velocity addition formula}
Turn off the moving grid and turn on the stationary grid. Make a new world line for an object moving with velocity $u$. This object will be moving in both the stationary and moving frames (since $u\neq v$). Set $u$ to such a value that it passes through one of the grid points.\footnote{There should be a folder that already does this} How is $u$ related to that grid point?

Switch back to the moving grid. Tune $u$ and $v$ such that the $u$ world line passes through a moving grid point. Make a note of the moving grid coordinates ($x',t'$). Given those moving grid coordinates, what does the moving observer find the velocity of the object to be ($u'$)? Make a note of the numerical values of all three velocities $u,v,u'$.

Repeat this process; try to find $u,v$ such that $|u-v|>c$. Make a note of $u,v,u'$ again in this case.

For a general $u$, transform the point $(x=u,t=1)$ into its coordinates in the moving frame ($x',t'$). A velocity is distance over time, so $u'=x'/t'$. With a little algebra you should get
$$
u'=\frac{u-v}{1-uv/c^2}
$$
Check that your two examples you wrote down fit this equation (plug in $u,v$).

\section{Report}
Some things to make sure you explain in your report:
\begin{enumerate}
    \item How are events, world lines, simultaneity, and light cones read off of a Minkowski diagram?
    \item What do the world lines of moving objects look like?
    \item How does the Lorentz transform give us the form of a moving grid in a Minkowski diagram and what does that look like? 
    \item Explain from the appearance of the moving grid: the second postulate (speed of light is constant for all observers) and the relativity of simultaneity.
    \item Explain using the two grids and a world line why the two frames observe different velocities for a single object moving relative to both.
    \item Explain how this leads to the velocity addition formula.
\end{enumerate}






\end{document}